\section{Appendix: Scope and Qualifications}

This study concerns the universal-calendar celebration of the Twenty-fourth
and Last Sunday after Pentecost on 22 November 2026 in the 1962 Roman
Missal. In this year's twenty-six Sunday slots after Pentecost, the
printed Twenty-fourth formulary is reserved for the last place.
The second-class Sunday has Gloria, Creed and the prescribed Trinity
Preface. Saint Cecilia's ordinary third-class feast supplies no
commemoration. No particular calendar or institutional additional prayer
is included.\footnote{\work{Rubricae generales} 18(b), 91, 109, 111(b);
1962 Missal, pp. 422--424 and the Saint Cecilia heading, p. 722.
The ICRSP Province of France Ordo, entry for 22 November 2026,
corroborates the temporal appointment within its institutional scope:
\url{https://icrspfrance.fr/ordo.php?d=2026-11-01}.}

All eleven appointed texts are included. The Latin follows the 1962
controller, with the independently image-checked Pustet 1862 antecedent.
Ligatures, stress accents and line divisions are normalized.
The Secret retains the target's longer conclusion cue; English short
conclusions remain the historical witness's own. The Trinity Preface
includes its complete thanksgiving, with separately sourced dialogue
and Sanctus junction. This is not a collation of the complete 1962
Ordinary.

The primary English for all seven scriptural appointments is the registered
Gutenberg Douay--Rheims--Challoner text. The appointed-text section quotes
its complete cited verses unchanged, with Jeremiah's nonconsecutive verses
kept separate. Each appointment then gives the complete historical liturgical
form separately from Lasance/Walsh's revised-1945 expression in a later
printing, not an identified 1937 impression. These parallel quotations
distinguish Bible wording from liturgical introductions, adaptations and
repetitions without conflating the two. The non-scriptural prayers,
Preface dialogue and Sanctus retain Lasance/Walsh throughout.
Neither source is an approved English edition of the 1962 Missal.
The Postcommunion's awkward rendering is retained and explained; the
Latin asks for cure of the fault. No new English prayer, conflated
Bible adaptation or expanded conclusion has been supplied.

The inherited Latin and historical English are excluded from the
project's licence. Their United States publication bases are the
registered public-domain antecedent and nonrenewal determinations.
Those findings do not clear the modern whole facsimile deliveries,
website wrappers or later inserts. No source scan is reproduced here.
Aquinas is explained from the checked Latin without reproducing the
restricted modern delivery. Bellarmine is paraphrased from the historical
abridged English transcription, without claiming a new print collation.
The NPNF English is historical translation, not a fresh Greek collation;
its Augustine Psalm text is abridged.

The two whole-Mass interpretations are editorial syntheses over the
checked explanations of particular passages. Lucien Fromage's continuation
of \work{The Liturgical Year} directly comments on the Mass's ten
non-Preface appointments; its older additional-oration discipline is
excluded. No direct patristic commentary on the three exact orations
or this exact Preface was established in the bounded research.
Doctrinal analogues are identified at their use. Bede's pertinent
paragraphs expound the Communion's immediate context, not a separately
quoted verse-24 lemma. Augustine's whole-Psalm explanations extend
beyond the verses appointed in the chants.

The generated chronology appendix distinguishes composition, narrated
episode, attribution and unresolved dates. Full source identities,
inspection extents, rights and search limits remain in
\texttt{propers/verified.md}, \texttt{research/scope.md} and the
leaf's source bindings. This is a study aid, not an official liturgical
book or an ecclesiastically approved commentary.

\section{References}

\begin{itemize}
\item \textbf{Missal and textual controls.}
\work{Missale Romanum}, Vatican typical edition, 1962, pp. 422--424,
nos. 1769--1778; Trinity Preface, p. 293, no. 1082;
\work{Rubricae generales} 18(b), 91, 109, 111(b).
Controlling facsimile:
\url{https://media.churchmusicassociation.org/pdf/missale62.pdf}.
\work{Missale Romanum} (Ratisbon: Pustet, 1862), pp. 359--360
and 251--252; original page-image archive, Internet Archive item
\texttt{bub\_gb\_E7sPAAAAIAAJ}.
Francis Xavier Lasance and Francis Augustine Walsh,
\work{The New Roman Missal}, revised-1945 expression, Benziger,
later c. 1955--56 printing: Last Sunday after Pentecost,
printed pp. 736--740; Ordinary dialogue, visible p. 772
(PDF 770); Trinity Preface and Sanctus, printed pp. 773--774.
The Latin misprint in Lasance's Gospel is not adopted.

\item \textbf{Scriptural English and biblical context.}
Douay--Rheims, Challoner revision, registered Project Gutenberg
no. 1581 per-book artifacts: primary English for the Introit's
Jer 29:11, 12, 14 and Ps 84(85):2; Epistle, Col 1:9--14;
Gradual, Ps 43(44):8--9; Alleluia and Offertory, Ps 129(130):1--2;
Gospel, Mt 24:15--35; Communion, Mk 11:24; also used for the
cited biblical contexts. This is a composite electronic Challoner
text, not one identified historical printing.
\url{https://www.gutenberg.org/ebooks/1581}.
Clementine Vulgate, registered
eBible \texttt{latVUC} per-book artifacts, for the textual
comparisons with the appointed Latin. The complete appointment
and comparison record is \texttt{propers/verified.md}.
Neither electronic Bible is treated as a witness to the
chants' liturgical redactions.

\item \textbf{Jerome.}
\work{In Ieremiam} V, on 29:10--14, in J.-P. Migne,
\work{Patrologia Latina} 24 (Paris, 1865), cols. 893--896;
facsimile PDF pp. 451--452. Original English exposition
of the Latin. Source:
\url{https://archive.org/details/sanctieusebiihie04jero}.

\item \textbf{Augustine.}
\work{Enarrationes in Psalmos}, Psalm 43, §§4--8, 13;
Psalm 84, §§1--13; Psalm 129, §§1--8. Historical abridged
English in \work{Nicene and Post-Nicene Fathers}, first series,
vol. VIII, ed. Philip Schaff, CCEL text; its English headings
are Psalms XLIV, LXXXV and CXXX.
\work{De Trinitate} III, preface §3; III.4.10;
III.10.21 and opening of III.11, historical English
translation by Arthur West Haddan (1887), registered
Book III text. Used as doctrinal illumination of the
Trinity and sacramental health, not direct commentary
on the Mass's prayers.

\item \textbf{John Chrysostom.}
\work{Homilies on Colossians} 2, on Col 1:9--14,
historical English translation attributed in the volume
to J. Ashworth, NPNF first series, vol. XIII,
CCEL text. \work{Homilies on Matthew} 75--77,
the checked exposition of Mt 24:15--35, NPNF first
series, vol. X, CCEL text. The study does not adopt
the sources' collective anti-Jewish punishment rhetoric.

\item \textbf{Thomas Aquinas.}
\work{Super Colossenses} 1, lectiones 3--4,
native anchors 87848--87849: the appointed prayer,
worthiness and rescue, followed by the Son's
divinity and relation to creation.
\url{https://www.corpusthomisticum.org/ccl.html}.
\work{Super Matthaeum} 24, lectio 3,
native anchor 87397, on the transition to the
final coming, deception, public manifestation
and the believing generation.
\url{https://www.corpusthomisticum.org/cml21.html}.
Original English exposition from the registered
restricted Latin deliveries.

\item \textbf{Bede.}
\work{In Marci evangelium expositio} III,
on Mk 11:22--23 and 25, paragraphs beginning
\textit{Et respondens Jesus} and
\textit{Et cum stabitis ad orandum},
registered Latin Migne-derived Wikisource
transcription. Immediate-context exegesis
of faith, the mountain, charity and forgiveness.

\item \textbf{Robert Bellarmine.}
\work{Commentary on the Book of Psalms},
trans. John O'Sullivan (Dublin: James Duffy,
1866), abridged historical English transcription,
eCatholic2000 delivery: Psalm 43, corresponding
explanation of rescue and ensuing lament,
especially native vv. 7--9; Psalm 129,
explanations of vv. 1--2 and 4--5.
Used by attributed paraphrase.

\item \textbf{Lucien Fromage.}
Continuation of Prosper Guéranger's
\work{The Liturgical Year}, trans.
Laurence Shepherd, vol. XI (Dublin:
James Duffy, 1883), ``Twenty-fourth
and Last Sunday after Pentecost,''
Mass commentary, pp. 503--504 and
512--514; in particular the Communion
and concluding explanation at p. 513.
This continuation is not attributed
to Guéranger personally.
\url{https://archive.org/details/liturgicalyear11gu}.
\end{itemize}
